\documentclass[11pt]{article}

\usepackage[margin=1in]{geometry}
\usepackage[T1]{fontenc}
\usepackage[utf8]{inputenc}
\usepackage{lmodern}
\usepackage{microtype}
\usepackage{amsmath}
\usepackage{array}
\usepackage{booktabs}
\usepackage{tabularx}
\usepackage{enumitem}
\usepackage[table]{xcolor}
\usepackage{hyperref}
\usepackage{fancyhdr}

\definecolor{ReportBlue}{HTML}{1F4E79}
\definecolor{ReportTeal}{HTML}{1E847F}
\definecolor{ReportGray}{HTML}{F4F6F8}
\definecolor{ReportLine}{HTML}{D7DEE3}
\definecolor{ReportText}{HTML}{1F2933}

\hypersetup{
  colorlinks=true,
  linkcolor=ReportBlue,
  urlcolor=ReportTeal,
  pdftitle={Skill-JEPA-WM Implementation Report},
  pdfauthor={Codex}
}

\pagestyle{fancy}
\fancyhf{}
\fancyhead[L]{\textcolor{ReportBlue}{Skill-JEPA-WM Report}}
\fancyhead[R]{2026-04-11}
\fancyfoot[C]{\thepage}
\setlength{\headheight}{14pt}
\setlist[itemize]{leftmargin=1.5em, itemsep=0.35em, topsep=0.4em}

\newcolumntype{Y}{>{\raggedright\arraybackslash}X}

\begin{document}

\begin{center}
{\LARGE\bfseries \textcolor{ReportBlue}{Skill-JEPA-WM Implementation and Experiment Report}}\\[0.5em]
{\large Push-T Debug Track in \texttt{jepa-wms}}\\[0.4em]
{\normalsize Branch: \texttt{skill-jepa-wm}}
\end{center}

\vspace{0.75em}

\noindent
\fcolorbox{ReportLine}{ReportGray}{
\begin{minipage}{0.97\linewidth}
\textbf{Executive Summary.}
I implemented the Skill-JEPA-WM pipeline inside \texttt{jepa-wms}, corrected several evaluation inconsistencies, and reran the debug experiments on the local RTX 4070 8GB.
The final corrected protocol shows that the \textbf{10\%-label hierarchical model outperforms both the flat variant and the labeled-only baseline} on Push-T debug.
Under the best stable setting (\texttt{K=4}, \texttt{execute\_actions\_per\_plan=4}), the hierarchical controller reaches \textbf{25\% success} with mean pose distance \textbf{254.61}, while the flat controller remains at \textbf{0\% success} with pose distance \textbf{459.92}, and the labeled-only baseline remains at \textbf{0\% success} with pose distance \textbf{793.81}.
\end{minipage}
}

\section*{1. Goal}

The target was to build and validate a hierarchical JEPA-based world model with the following structure:

\begin{itemize}
\item a frozen V-JEPA2 encoder and offline feature cache,
\item a passive high-level latent skill model trained on mostly action-free data,
\item a low-level action-conditioned JEPA world model trained on a small labeled subset,
\item hierarchical planning with skill-level MPC above primitive-action MPC.
\end{itemize}

The concrete success criteria from the design spec were:

\begin{itemize}
\item no latent collapse,
\item the hierarchical planner beats the flat primitive-only planner,
\item the 10\%-label setup beats the labeled-only low-level baseline at the same label budget.
\end{itemize}

\section*{2. What Was Implemented}

The implementation added a new \texttt{skill\_jepa/} package and the experiment plumbing required to train and evaluate it inside \texttt{jepa-wms}.
The main work items were:

\begin{itemize}
\item frozen V-JEPA2 feature extraction and offline HDF5 caching,
\item state projection to global and spatial JEPA latents,
\item passive skill inference, skill world model, and skill prior,
\item low-level action chunk encoder and action-conditioned JEPA world model,
\item hierarchical planners and analysis scripts,
\item Push-T debug configs, ablations, and experiment outputs.
\end{itemize}

\section*{3. Incompletenesses Found During Review}

The main value of the second pass was not more code volume. It was correcting invalid or incomplete experiment logic.

\subsection*{3.1 Representation Cache Mismatch}

The initial cache path encoded duplicated single frames as a two-frame video input.
That weakened temporal signal and made start/goal latent distances unrealistically small even when simulator state distance was large.

\textbf{Fix.}
I changed the cache builder and encoder path to use causal two-frame clips.

\subsection*{3.2 Online Evaluation Mismatch}

The early online evaluator had several inconsistencies:

\begin{itemize}
\item it sampled goals farther in the future than the controller was allowed to act,
\item it mixed velocity terms into the reported task distance,
\item it encoded goals from a fresh environment render instead of using the cached goal latents already used for training,
\item the planner comparison was order-sensitive because the stochastic planner state was not reset per episode.
\end{itemize}

\textbf{Fix.}
I corrected all four issues:

\begin{itemize}
\item capped goal sampling to the control horizon,
\item reported pose distance on the task state,
\item used cache-consistent start and goal latents,
\item reseeded each evaluation episode deterministically.
\end{itemize}

\subsection*{3.3 Control Execution Policy}

The original online evaluator always executed one primitive action before replanning.
The spec explicitly allows executing the first \(M\) primitive actions and then replanning.

\textbf{Fix.}
I added \texttt{execute\_actions\_per\_plan} to the evaluator and found that \texttt{M=4} improves the best stable debug run.

\section*{4. Experiment Protocol}

\subsection*{4.1 Environment}

\begin{itemize}
\item Local machine: Windows workspace, local RTX 4070 8GB
\item Encoder: \texttt{facebook/vjepa2-vitl-fpc64-256}
\item Dataset: readable Push-T debug backup
\item Training mode: local only, no remote GPU rental
\end{itemize}

\subsection*{4.2 Final Main Debug Setting}

\begin{itemize}
\item chunk size \(K = 4\),
\item control horizon \(= 32\),
\item goal gap \(= 24\),
\item execute first \(M = 4\) primitive actions per plan,
\item corrected cache-consistent online evaluation.
\end{itemize}

\subsection*{4.3 Sanity Check}

The corrected online task set is reachable under oracle replay of cached expert actions:

\begin{center}
\begin{tabular}{lr}
\toprule
Metric & Value \\
\midrule
Oracle replay success rate & 1.00 \\
Oracle replay mean pose distance & 2.55 \\
\bottomrule
\end{tabular}
\end{center}

This matters because it shows that failures in the learned controller are model or planning failures, not an impossible evaluation setup.

\section*{5. Final Results}

\subsection*{5.1 Training Stability}

\begin{center}
\begin{tabularx}{0.96\linewidth}{lrrrr}
\toprule
\rowcolor{ReportBlue!10}
Stage & Skill End & Skill Roll & Skill Std & Posterior Var \\
\midrule
Passive (epoch 3) & 0.0778 & 0.1473 & 0.0350 & 2.9512 \\
Joint (epoch 3) & 0.0530 & 0.0847 & 0.0295 & 7.2426 \\
\bottomrule
\end{tabularx}
\end{center}

The passive and joint runs show no latent collapse under the criterion used in the design spec.
The skill standard deviation remains positive, and rollout losses decrease substantially.

\subsection*{5.2 Final Online Comparison}

\begin{center}
\small
\begin{tabular}{llllll}
\toprule
\rowcolor{ReportBlue!10}
Run & Success & Pose Dist. & Planner Latency & Better Rate & Comment \\
\midrule
Joint hierarchical (\(K=4, M=4\)) & 0.25 & 254.61 & 0.212 s & 0.75 & Best final run \\
Joint flat (\(K=4, M=4\)) & 0.00 & 459.92 & 0.525 s & -- & Same checkpoint, flat planner \\
Labeled-only flat baseline & 0.00 & 793.81 & 0.527 s & -- & Same label budget \\
Joint hierarchical (\(K=2\) ablation) & 0.00 & 615.69 & 0.282 s & 0.375 & Weaker than final \(K=4\) \\
\bottomrule
\end{tabular}
\end{center}

\subsection*{5.3 Interpretation}

The corrected final result supports three claims:

\begin{itemize}
\item The hierarchical planner beats the flat primitive-only variant on the same learned model.
\item The 10\%-label hierarchical method beats the labeled-only low-level baseline at the same label budget.
\item The implementation is complete enough to satisfy the first complete-run criteria from the spec.
\end{itemize}

The main remaining limitation is absolute control quality.
The learned controller is still well below the oracle replay baseline on this small debug backup.

\section*{6. Summary of Concrete Fixes}

\begin{center}
\small
\begin{tabular}{ll}
\toprule
\rowcolor{ReportBlue!10}
Area & Fix \\
\midrule
Feature cache & Replaced duplicated single-frame inputs with causal two-frame clips \\
Goal construction & Restricted online goals to the executable control budget \\
State metric & Reported pose distance instead of mixing in velocity terms \\
Goal encoding & Used cache-consistent start and goal latents for online evaluation \\
Planner comparison & Reseeded each evaluation episode deterministically \\
Execution policy & Added execution of the first \(M\) primitive actions before replanning \\
Push-T runtime & Patched local gymnasium and pymunk compatibility issues \\
\bottomrule
\end{tabular}
\end{center}

\section*{7. Conclusion}

This work completed the Skill-JEPA-WM debug pipeline, corrected the evaluation protocol, and produced a final stable result that satisfies the initial milestone.
The best current setting is the corrected \texttt{K=4} model with \texttt{execute\_actions\_per\_plan=4}.
Its online Push-T debug result is:

\begin{center}
\fcolorbox{ReportBlue}{ReportBlue!6}{
\begin{minipage}{0.9\linewidth}
\textbf{Final headline result.}
Hierarchical control reached \textbf{25\% success} and mean pose distance \textbf{254.61}, while the matched flat controller stayed at \textbf{0\% success} and the labeled-only baseline also stayed at \textbf{0\% success}.
\end{minipage}
}
\end{center}

The next step is not more plumbing.
The next step is scaling the winning protocol to a larger Push-T or DROID subset and improving control quality against the oracle replay ceiling.

\end{document}
